\documentclass[12pt]{article}
\usepackage[utf8]{inputenc}
\usepackage[T1]{fontenc}
\usepackage{hyperref}
\title{Maya, Chan Santa Cruz Language Resources for K-12 Curriculum Development}
\author{Compiled by Claude Sonnet 5.5 (claude-sonnet-5-5) for Andr\'as Kornai}
\date{}
\begin{document}
\maketitle

\section{Introduction}
Maya, Chan Santa Cruz (ISO 639-3 yus, Glottocode maya1286) is the Maya variety of the Chan Santa Cruz area, now Felipe Carrillo Puerto, Quintana Roo, Mexico. The family given in the prompt (``Bookkeeping'') is a catalogue placeholder rather than a real family name. The language is closely tied to Yucatec Maya. This report rests on a single search, and nothing was found that is specific to this variety.

\section{Primary Educational Resources}
A Harvard DRCLAS article describes a bilingual (Maya/Spanish) education degree programme at the Universidad Bienestar Benito Ju\'arez in Felipe Carrillo Puerto: \url{https://revista.drclas.harvard.edu/tearing-down-the-walls-of-education-the-maya-struggle-against-colonialism}

\section{Linguistic and Reference Materials}
\begin{itemize}
\item A thesis on Maya language use from CIESAS: \url{https://ciesas.repositorioinstitucional.mx/jspui/bitstream/1015/1667/1/TE%20E.E.%202023%20Lendy%20Rosario%20Ek%20Ek.pdf}
\item El Colegio de M\'exico grammar-archive volume on Yucatec Maya of Xhazil Sur, which mentions the Diccionario de Motul and the Academia de la Lengua y Cultura Mayas de Quintana Roo in Felipe Carrillo Puerto: \url{https://cell.colmex.mx/ast/blobs/redirect/eyJfcmFpbHMiOnsibWVzc2FnZSI6IkJBaHBhdz09IiwiZXhwIjpudWxsLCJwdXIiOiJibG9iX2lkIn19--c85ed0c352a05835cc275764002e7dbab643bdfa/32_maya-yucateco-de-xhazil-sur.pdf}
\end{itemize}

\section{Religious and Community Resources}
The Academia de la Lengua y Cultura Mayas de Quintana Roo, named in the Colmex volume above, is a community institution. No website was found.

\section{Digital and Audio Resources}
An AILLA (University of Texas) archive record titled ``Chan Santa Cruz'': \url{https://islandora-ailla.lib.utexas.edu/islandora/object/ailla%3A126433}
Its contents were not verified.

\section{Curriculum Development Context}
Yucatec Maya materials are likely to be the main base, with local adaptation. No variety-specific textbooks or dictionaries were found in this search.
\end{document}